\documentclass[10pt]{beamer}
\usetheme{metropolis}
\usepackage{graphicx}
\usepackage{listings}
\usepackage{booktabs}
\usepackage{xcolor}
\usepackage{hyperref}

\definecolor{codebg}{HTML}{F5F5F5}
\definecolor{codegreen}{HTML}{2E7D32}
\definecolor{codegray}{HTML}{757575}
\definecolor{codeblue}{HTML}{1565C0}

\lstdefinestyle{rstyle}{
  language=R, backgroundcolor=\color{codebg},
  basicstyle=\ttfamily\scriptsize,
  keywordstyle=\color{codeblue}\bfseries,
  stringstyle=\color{codegreen},
  commentstyle=\color{codegray}\itshape,
  breaklines=true, frame=single,
  rulecolor=\color{codegray!30}, showstringspaces=false, tabsize=2,
  morekeywords={library,ggplot,aes,geom_point,geom_tile,geom_col,
    scale_fill_viridis_c,scale_color_brewer,scale_fill_manual,
    theme_minimal,labs,coord_flip}
}
\lstdefinestyle{pythonstyle}{
  language=Python, backgroundcolor=\color{codebg},
  basicstyle=\ttfamily\scriptsize,
  keywordstyle=\color{codeblue}\bfseries,
  stringstyle=\color{codegreen},
  commentstyle=\color{codegray}\itshape,
  breaklines=true, frame=single,
  rulecolor=\color{codegray!30}, showstringspaces=false, tabsize=4,
  morekeywords={import,from,as,pd,plt,sns,np,fig,ax}
}

\title{Color Theory \& Accessibility}
\subtitle{Workshop 1 --- Module 4}
\author{Prof.Asoc.\ Endri Raco}
\institute{Data Visualization for Data Scientists \\ UNYT Tirana}
\date{2025/26}

\begin{document}

\begin{frame}
  \titlepage
\end{frame}

% ================================================================
\begin{frame}{Module Roadmap}
  \begin{enumerate}
    \item The three dimensions of colour: hue, saturation, lightness
    \item Palette typology: sequential, diverging, qualitative
    \item Perceptual uniformity and the viridis family
    \item Colour vision deficiency and accessible design
    \item WCAG contrast standards and the grey-plus-accent strategy
  \end{enumerate}
  \begin{exampleblock}{Learning Outcome}
    You will be able to select, justify, and implement colour palettes
    that are perceptually accurate, colourblind-safe, and WCAG-compliant.
  \end{exampleblock}
\end{frame}

% ================================================================
\section{The Three Dimensions of Colour}
% ================================================================

\begin{frame}{Hue, Saturation, Lightness}
  \begin{center}
    \includegraphics[width=0.95\textwidth]{figures_m04/hsl}
  \end{center}
  \begin{itemize}
    \item \textbf{Hue}: the colour name (red, blue, green) --- categorical channel
    \item \textbf{Saturation}: vividness vs.\ grey --- ordinal / emphasis channel
    \item \textbf{Lightness}: dark to bright --- sequential / quantitative channel
  \end{itemize}
\end{frame}

% ================================================================
\begin{frame}{Mapping Data to Colour Dimensions}
  \centering
  \begin{tabular}{lll}
    \toprule
    \textbf{Data Type} & \textbf{Use} & \textbf{Avoid} \\
    \midrule
    Nominal (categories) & Hue differences          & Lightness gradient \\
    Ordinal (rank)       & Lightness / saturation   & Hue rainbow \\
    Quantitative (value) & Lightness gradient       & Discrete hues \\
    Diverging ($\pm$)    & Two-hue + neutral centre & Single hue \\
    \bottomrule
  \end{tabular}

  \vspace{8pt}
  \begin{alertblock}{Pro-Tip}
    The single most common colour error is using a \textbf{rainbow} (hue
    variation) for quantitative data. Rainbow palettes create false
    perceptual boundaries where the data are smooth.
  \end{alertblock}
\end{frame}

% ================================================================
\section{Palette Types}
% ================================================================

\begin{frame}{Sequential, Diverging, Qualitative}
  \begin{center}
    \includegraphics[width=0.9\textwidth]{figures_m04/palette_types}
  \end{center}
  \begin{exampleblock}{Data Insight}
    Choosing the wrong palette type is a semantic error: a diverging
    palette implies a meaningful midpoint; a qualitative palette implies
    no ordering. Match the palette structure to the data structure.
  \end{exampleblock}
\end{frame}

% ================================================================
\begin{frame}{ColorBrewer: Empirically Tested Palettes}
  \begin{center}
    \includegraphics[width=0.85\textwidth]{figures_m04/colorbrewer}
  \end{center}
  \begin{itemize}
    \item Brewer (2003) tested palettes for cartographic legibility
    \item Three families: sequential, diverging, qualitative
    \item Each family rated for print, photocopy, LCD, and colourblind safety
    \item Interactive tool: \texttt{https://colorbrewer2.org}
  \end{itemize}
\end{frame}

% ================================================================
\begin{frame}{The Palette Decision Tree}
  \begin{center}
    \includegraphics[width=0.85\textwidth]{figures_m04/decision_tree}
  \end{center}
\end{frame}

% ================================================================
\section{Perceptual Uniformity}
% ================================================================

\begin{frame}{Why Rainbow Fails}
  \begin{center}
    \includegraphics[width=0.92\textwidth]{figures_m04/rainbow_vs_viridis}
  \end{center}
  \begin{alertblock}{Pro-Tip}
    Rainbow (jet) creates false contour bands at yellow and cyan where
    the human visual system perceives abrupt lightness jumps. Viridis
    changes lightness monotonically, so equal data steps produce equal
    perceptual steps.
  \end{alertblock}
\end{frame}

% ================================================================
\begin{frame}{The Viridis Family}
  \begin{center}
    \includegraphics[width=0.85\textwidth]{figures_m04/viridis_family}
  \end{center}
  \begin{itemize}
    \item Designed by van der Walt \& Smith (2015) for matplotlib
    \item Perceptually uniform in both colour and greyscale
    \item Safe for deuteranopia, protanopia, and tritanopia
    \item Available in R (\texttt{viridis}) and Python (\texttt{plt.cm.viridis})
  \end{itemize}
\end{frame}

% ================================================================
\begin{frame}[fragile]{Sequential Palette in R}
  \begin{columns}[T]
    \begin{column}{0.52\textwidth}
\begin{lstlisting}[style=rstyle]
library(ggplot2)
library(viridis)

# Heatmap with perceptually
# uniform sequential palette
ggplot(df, aes(x = week,
               y = day,
               fill = value)) +
  geom_tile(color = "white",
            linewidth = 0.3) +
  scale_fill_viridis_c(
    option = "D",     # viridis
    direction = 1) +
  theme_minimal(base_size = 9) +
  labs(title = "Weekly Heatmap",
       fill = "Value")

# Alternatives:
# option = "A" (magma)
# option = "B" (inferno)
# option = "C" (plasma)
\end{lstlisting}
    \end{column}
    \begin{column}{0.45\textwidth}
      \includegraphics[width=\textwidth]{figures_m04/heatmap_r}
    \end{column}
  \end{columns}
\end{frame}

% ================================================================
\begin{frame}[fragile]{Sequential Palette in Python}
  \begin{columns}[T]
    \begin{column}{0.52\textwidth}
\begin{lstlisting}[style=pythonstyle]
import matplotlib.pyplot as plt
import numpy as np

data = np.random.uniform(
    10, 90, (6, 8))

fig, ax = plt.subplots()
im = ax.imshow(
    data,
    cmap="viridis",  # default
    aspect="auto")

plt.colorbar(im, ax=ax,
    shrink=0.8, label="Value")

ax.set_title("Weekly Heatmap")
plt.tight_layout()

# Alternatives:
# cmap="magma"
# cmap="plasma"
# cmap="inferno"
# cmap="cividis" (colourblind)
\end{lstlisting}
    \end{column}
    \begin{column}{0.45\textwidth}
      \includegraphics[width=\textwidth]{figures_m04/heatmap_py}
    \end{column}
  \end{columns}
\end{frame}

% ================================================================
\begin{frame}[fragile]{Diverging Palette in R}
  \begin{columns}[T]
    \begin{column}{0.52\textwidth}
\begin{lstlisting}[style=rstyle]
# Diverging: centred on zero
# Two hues + white midpoint
ggplot(df, aes(x = x, y = y,
               color = deviation)) +
  geom_point(size = 2.5) +
  scale_color_gradient2(
    low = "#2166AC",    # blue
    mid = "white",      # neutral
    high = "#B2182B",   # red
    midpoint = 0,
    limits = c(-1, 1)) +
  theme_minimal() +
  labs(color = "Deviation")

# ColorBrewer alternative:
# scale_color_distiller(
#   palette = "RdBu",
#   direction = -1)
\end{lstlisting}
    \end{column}
    \begin{column}{0.45\textwidth}
      \includegraphics[width=\textwidth]{figures_m04/diverging_map}

      \vspace{4pt}
      {\small The midpoint (white) must align
      with the meaningful centre of the data
      (usually zero or the mean).}
    \end{column}
  \end{columns}
\end{frame}

% ================================================================
\begin{frame}[fragile]{Diverging Palette in Python}
  \begin{columns}[T]
    \begin{column}{0.52\textwidth}
\begin{lstlisting}[style=pythonstyle]
import matplotlib.pyplot as plt

fig, ax = plt.subplots()

sc = ax.scatter(
    x, y, c=deviation,
    cmap="RdBu_r",      # diverging
    vmin=-1, vmax=1,     # centre=0
    s=60,
    edgecolors="white",
    linewidth=0.5)

plt.colorbar(sc, ax=ax,
    label="Deviation from Mean")

# TwoSlopeNorm for off-centre
from matplotlib.colors import (
    TwoSlopeNorm)
norm = TwoSlopeNorm(
    vmin=-2, vcenter=0, vmax=5)
# Pass norm=norm to scatter()
\end{lstlisting}
    \end{column}
    \begin{column}{0.45\textwidth}
      \textbf{Key parameters:}

      \begin{itemize}
        \item \texttt{vmin}, \texttt{vmax}: symmetric
              around zero
        \item \texttt{TwoSlopeNorm}: for asymmetric
              ranges (e.g., $-2$ to $+5$)
        \item Suffix \texttt{\_r}: reverses the
              colourmap direction
      \end{itemize}
    \end{column}
  \end{columns}
\end{frame}

% ================================================================
\section{Colour Vision Deficiency}
% ================================================================

\begin{frame}{8\% of Males Have Colour Vision Deficiency}
  \begin{center}
    \includegraphics[width=0.92\textwidth]{figures_m04/colorblind_sim}
  \end{center}
  \begin{exampleblock}{Data Insight}
    Deuteranopia (red--green deficiency) affects $\sim$8\% of males and
    $\sim$0.5\% of females. Red and green bars become nearly
    indistinguishable. Always simulate CVD before publishing.
  \end{exampleblock}
\end{frame}

% ================================================================
\begin{frame}{Accessible Design Strategies}
  \begin{enumerate}
    \item Use \textbf{viridis} or \textbf{cividis} for sequential data ---
          both are colourblind-safe by construction
    \item Supplement colour with a \textbf{second channel}: shape, pattern,
          direct label, or line style
    \item Simulate CVD using tools:
          \begin{itemize}
            \item R: \texttt{colorspace::deutan()}, \texttt{protan()}, \texttt{tritan()}
            \item Python: \texttt{colorspacious} package
            \item Web: \texttt{https://www.color-blindness.com/coblis/}
          \end{itemize}
    \item Limit qualitative palettes to $\leq$ 7 hues --- CVD simulation
          degrades rapidly beyond this
  \end{enumerate}
\end{frame}

% ================================================================
\begin{frame}[fragile]{CVD Simulation in R}
  \begin{columns}[T]
    \begin{column}{0.52\textwidth}
\begin{lstlisting}[style=rstyle]
library(colorspace)

# Create a plot
p <- ggplot(df, aes(x, y,
    color = category)) +
  geom_point(size = 3) +
  scale_color_manual(
    values = c("A" = "#E53935",
               "B" = "#2E7D32",
               "C" = "#1565C0"))

# Simulate deuteranopia
protan(p)   # protanopia
deutan(p)   # deuteranopia
tritan(p)   # tritanopia

# Better: use safe palettes
p + scale_color_viridis_d()
# or
p + scale_color_brewer(
    palette = "Set2")
\end{lstlisting}
    \end{column}
    \begin{column}{0.45\textwidth}
      \textbf{Workflow:}

      \begin{enumerate}
        \item Build chart with default palette
        \item Run \texttt{deutan(p)} to simulate
        \item If categories merge, switch to
              a colourblind-safe palette
        \item Add a redundant channel (shape)
              as insurance
      \end{enumerate}

      \vspace{4pt}
      {\small The \texttt{colorspace} package also
      provides \texttt{hcl\_palettes()} for
      systematic palette exploration.}
    \end{column}
  \end{columns}
\end{frame}

% ================================================================
\begin{frame}[fragile]{CVD Simulation in Python}
  \begin{columns}[T]
    \begin{column}{0.52\textwidth}
\begin{lstlisting}[style=pythonstyle]
# pip install colorspacious
from colorspacious import cspace_convert
import numpy as np

def simulate_cvd(rgb_array,
                 deficiency="deuteranomaly",
                 severity=100):
    """Simulate colour vision
    deficiency on an RGB image."""
    # Convert to CVD space
    cvd = cspace_convert(
        rgb_array,
        {"name": "sRGB1",
         "CVD": {
           "type": deficiency,
           "severity": severity}},
        "sRGB1")
    return np.clip(cvd, 0, 1)

# Alternative: use cividis cmap
# plt.scatter(x, y, c=z,
#     cmap="cividis")
# cividis designed for CVD safety
\end{lstlisting}
    \end{column}
    \begin{column}{0.45\textwidth}
      \textbf{Quick alternatives:}

      \begin{itemize}
        \item \texttt{cmap="cividis"} is optimised
              for deuteranopia and protanopia
        \item Seaborn's \texttt{palette="colorblind"}
              provides a 6-colour safe set
        \item Add markers: \texttt{style="category"}
              in \texttt{sns.scatterplot()}
      \end{itemize}
    \end{column}
  \end{columns}
\end{frame}

% ================================================================
\section{Contrast \& the Grey-Accent Strategy}
% ================================================================

\begin{frame}{WCAG 2.1 Contrast Requirements}
  \begin{center}
    \includegraphics[width=0.92\textwidth]{figures_m04/wcag_contrast}
  \end{center}
  \begin{itemize}
    \item \textbf{AA}: contrast ratio $\geq$ 4.5:1 for normal text,
          $\geq$ 3:1 for large text (18pt+)
    \item \textbf{AAA}: contrast ratio $\geq$ 7:1
    \item Check with: \texttt{https://webaim.org/resources/contrastchecker/}
  \end{itemize}
\end{frame}

% ================================================================
\begin{frame}{Grey + Accent: The Most Powerful Colour Strategy}
  \begin{center}
    \includegraphics[width=0.92\textwidth]{figures_m04/grey_accent}
  \end{center}
  \begin{exampleblock}{Data Insight}
    Colouring everything equally gives everything equal salience ---
    which means nothing stands out. Grey desaturates the background;
    a single accent hue creates instant pre-attentive pop-out for the
    focal data.
  \end{exampleblock}
\end{frame}

% ================================================================
\begin{frame}[fragile]{Grey + Accent in R}
  \begin{columns}[T]
    \begin{column}{0.52\textwidth}
\begin{lstlisting}[style=rstyle]
library(ggplot2)

# Highlight one category, grey rest
ggplot(df, aes(x = x, y = y,
    color = is_highlight)) +
  geom_point(aes(
    size = ifelse(is_highlight,
                  3, 1.5),
    alpha = ifelse(is_highlight,
                   0.9, 0.3))) +
  scale_color_manual(
    values = c(
      "TRUE" = "#E53935",
      "FALSE" = "#BBBBBB")) +
  scale_size_identity() +
  scale_alpha_identity() +
  theme_minimal() +
  theme(legend.position = "none") +
  labs(title = "Grey + Accent")
\end{lstlisting}
    \end{column}
    \begin{column}{0.45\textwidth}
      \textbf{Three-step recipe:}

      \begin{enumerate}
        \item Create a boolean column
              \texttt{is\_highlight}
        \item Map \texttt{TRUE} to accent,
              \texttt{FALSE} to grey
        \item Increase size and alpha for
              the highlighted subset
      \end{enumerate}

      \vspace{4pt}
      {\small This leverages both colour
      and size pre-attentive channels
      simultaneously.}
    \end{column}
  \end{columns}
\end{frame}

% ================================================================
\begin{frame}[fragile]{Grey + Accent in Python}
  \begin{columns}[T]
    \begin{column}{0.52\textwidth}
\begin{lstlisting}[style=pythonstyle]
import matplotlib.pyplot as plt

fig, ax = plt.subplots()

# Background points (grey)
mask = df["category"] != "target"
ax.scatter(
    df.loc[mask, "x"],
    df.loc[mask, "y"],
    c="#BBBBBB", s=20,
    alpha=0.3, edgecolors="none")

# Highlighted points (accent)
mask_hl = df["category"] == "target"
ax.scatter(
    df.loc[mask_hl, "x"],
    df.loc[mask_hl, "y"],
    c="#E53935", s=50,
    alpha=0.9,
    edgecolors="white",
    linewidth=0.5, zorder=5)

ax.set_title("Grey + Accent")
\end{lstlisting}
    \end{column}
    \begin{column}{0.45\textwidth}
      \textbf{Python pattern:}

      \begin{enumerate}
        \item Plot background first (grey,
              small, low alpha)
        \item Plot highlight second (accent,
              large, high alpha, \texttt{zorder=5})
        \item No legend needed --- the colour
              difference is self-explanatory
      \end{enumerate}
    \end{column}
  \end{columns}
\end{frame}

% ================================================================
\begin{frame}{Colour Do's and Don'ts}
  \begin{center}
    \includegraphics[width=0.92\textwidth]{figures_m04/dos_donts}
  \end{center}
\end{frame}

% ================================================================
\section{Synthesis}
% ================================================================

\begin{frame}{Key Takeaways}
  \begin{enumerate}
    \item \textbf{Hue} for categories, \textbf{lightness} for quantities,
          \textbf{saturation} for emphasis --- never conflate the three.
    \item \textbf{Sequential} for ordered data, \textbf{diverging} for
          data with a meaningful midpoint, \textbf{qualitative} for
          nominal categories ($\leq$ 7).
    \item \textbf{Never use rainbow}. Use \texttt{viridis}, \texttt{magma},
          or \texttt{cividis} for perceptual uniformity.
    \item \textbf{Simulate CVD} before publishing (\texttt{colorspace} in R,
          \texttt{colorspacious} in Python).
    \item \textbf{Grey + accent} is the single most effective colour strategy
          for explanatory visualization.
  \end{enumerate}
\end{frame}

% ================================================================
\begin{frame}{Essential Readings for This Module}
  \begin{itemize}
    \item Brewer, C.~A. (2003). ``A Transition in Improving Maps: The
          ColorBrewer Example.'' \emph{Cartography and Geographic
          Information Science}, 30(2).
    \item van der Walt, S. \& Smith, N. (2015). ``A Better Default
          Colormap for Matplotlib.'' SciPy 2015.
    \item Crameri, F., Shephard, G.~E. \& Heron, P.~J. (2020). ``The
          Misuse of Colour in Science Communication.''
          \emph{Nature Communications}, 11, 5444.
    \item Wilke, C.~O. (2019). \emph{Fundamentals of Data Visualization},
          Chapter 4 (Color Scales).
  \end{itemize}

  \vspace{6pt}
  \textbf{Next Module}: Typography, Layout \& Composition (W01-M05)
\end{frame}

\end{document}
